\chapter{Introduction}

Portable Document Format (PDF) is the dominant format for sharing and archiving
documents across law, finance, healthcare, academia and government. This is
largely because a PDF looks identical on every device and cannot be accidentally
edited. But that same property creates a serious problem for automated
processing. PDF stores
text as individual characters placed at specific coordinates on a page. There
is no native concept of paragraphs, headings, or lists. Any tool that tries to extract content from
a PDF has to reconstruct that structure from scratch, using only positional
and typographic signals, which is unreliable \autocite{whitington2012}.

Three main approaches exist for PDF-to-Markdown conversion. Text-based parsers
such as PyMuPDF and pdfplumber \autocite{meuschke2023} read the raw text layer of
a native PDF and recover the characters, but discard layout structure in the
process. Optical Character Recognition (OCR) tools such as Tesseract \autocite{smith2007} treat the page as a flat
image and recognize characters pixel by pixel. This works for scanned documents,
but the output is plain text with no structure. A newer generation of
AI-based systems, including Marker, Microsoft MarkItDown, LlamaParse, and
multimodal models such as Mistral OCR \autocite{omnidocbench2025}, uses deep learning
to recover more structure and handles tables, formulas, and multi-column layouts
more reliably. Many pipelines combine all three stages: scanned pages go through
OCR first, then a layout-aware conversion tool, then regex-based post-processing.

But all of these share one problem: they apply the same strategy to every page
of a document. A tool designed for text-heavy papers uses that same approach on
a page that contains only a scanned diagram or a dense formula. Real-world
documents mix plain text, tables, images, source code, and equations, sometimes
on the same page, and no single strategy handles all of them well. Ingemarsson
and Persson \autocite{ingemarsson2024} tested this directly: rule-based tools scored
above 86\% on plain-text PDFs but dropped to 0\% on image-only pages. No single
method came out on top across all content types. And to the best of the author's
knowledge, no widely used open-source tool picks a different strategy per page
based on what that page actually contains.

This thesis builds and evaluates a system that does exactly that. Instead of
applying one fixed strategy to the whole document, the system classifies each
page as plain text, image-heavy, formula-heavy, table-heavy, mixed content, or
empty, and routes it to the extraction method best suited to that type. The system runs locally using LM Studio as the Large Language Model (LLM)
runtime, so there is no dependency on external cloud services and results are
fully reproducible. Converting a PDF to Markdown is more than a format change ---
it means recovering the document's structure: headings, lists, tables, links,
code blocks, and hierarchy, from a format that encodes visual appearance but not
meaning. This chapter explains the problem, defines the research questions, and
outlines the scope of the work.

% ----------------------------------------------------------------
\section{Problem Statement}
% ----------------------------------------------------------------

The problem this thesis addresses has two parts: the PDF format itself hides
structural information, and existing tools cannot adjust their extraction
approach based on what a given page actually contains.

\subsection{Difficulty of Extracting Structured Content from PDFs}

PDF was not designed to store meaning. It was designed to reproduce the
visual appearance of a printed page. The format has no native representation
of paragraphs, headings, list items, or table cells \autocite{whitington2012}.
Any structure a reader sees has to be reconstructed by software from
low-level positional and typographic signals, which is fragile and often
wrong \autocite{zhu2022}. A single PDF can mix plain text, embedded images,
mathematical formulas, multi-column layouts, and complex tables, sometimes
on the same page. No single extraction strategy handles all of these
correctly. Meuschke et al.\ \autocite{meuschke2023} evaluated mainstream parsing
tools across multiple document domains and found that even the top performers
failed on scientific and patent documents.

\subsection{Limited OCR Accuracy for Complex Layouts}

OCR is the standard approach for scanned PDFs, where text exists only as
pixels. But OCR accuracy drops significantly on complex layouts. Subramani
et al.\ \autocite{subramani2020} describe OCR as a two-stage process: text
detection, which locates where text appears on the page, and text
transcription, which recognises characters within each detected region.
Layout analysis must happen before both stages: the tool first needs to
identify content regions (text blocks, figures, tables, equations) before
any recognition can begin. Blecher et al.\ \autocite{blecher2023} and Subramani
et al.\ \autocite{subramani2020} explain the mechanical limitation: traditional
OCR engines are built on a detect-first, recognise-second pattern where any
error in layout segmentation propagates directly into text recognition,
whereas modern unified visual models generate structured markup from full-page
screenshots without a separate layout step. Tools like
Tesseract~\autocite{smith2007} skip this step and treat the whole page as a
single flat image, which makes it hard to assign the right semantic role to
each detected region. Zhu et al.\ \autocite{zhu2022} measured the effect
directly: without layout segmentation, normalized edit distance was 0.63;
with a layout segmentation step added beforehand, it dropped to 0.38 (lower
is better). That is a large gain from one preprocessing step alone.

Multi-column layouts are a major source of errors. OCR engines often merge
text across columns or get the reading order wrong, and Zhu et
al.\ \autocite{zhu2022} show that preserving correct reading order in
multi-column documents requires a dedicated segmentation stage before any
character recognition is attempted. Even current models struggle with dense
multi-column text, rotated text, tables, and mathematical notation
\autocite{subramani2020, meuschke2023}. Blecher et al.\ \autocite{blecher2023}
note that the PDF format causes ``a loss of semantic information,
particularly for mathematical expressions,'' and Darvishy~\autocite{darvishy2018}
points out that no existing tool can interpret or describe mathematical
formulas and scientific graphs. The core issue is structural, not just
visual: what a page contains determines which extraction approach has any
chance of working.

\subsection{Need for Conversion to Markdown for Accessibility and Version Control}

Markdown is a lightweight markup language created by John Gruber in 2004.
It is readable as plain text and converts cleanly to HTML and other formats
\autocite{cone2019}.

Converting PDFs to Markdown matters for accessibility. Darvishy
\autocite{darvishy2018} explains the problem: PDF stores pages as rendering
instructions, not as structured content. Screen readers and other assistive
technologies need semantic tags and a defined reading order to process a
document correctly. Tagged PDF addresses this, but most PDFs in practice are
untagged, and Darvishy found that even when tags are present, correct reading
order and structure still require manual verification. For scientific documents
the gap is larger: no existing PDF accessibility tool can interpret
mathematical formulas or scientific graphs, which leaves that content
unreachable for screen reader users \autocite{darvishy2018}. Markdown marks structure directly in the text: headings with \texttt{\#},
lists with \texttt{-}, formulas with \texttt{\$}. Assistive tools can
process it without guessing. The GitHub Flavored Markdown extension adds table syntax and
\LaTeX-style math notation, which covers most of what academic documents
require \autocite{duan2025}.

Markdown's other advantages follow from being plain text. It works with
version control systems like Git, where every change is tracked at the
character level, which PDF and binary formats cannot offer. It is easy to
parse programmatically and stays readable without any specific software.
Khan et al.\ \autocite{khan2024} show a concrete consequence for AI pipelines:
the quality of text extraction sets a ceiling on retrieval precision in RAG
systems. Poorly structured text produces low-quality embeddings, and retrieval
fails even when the relevant content is present. Markdown's structure prevents
that degradation.

% ----------------------------------------------------------------
\section{Motivation}
% ----------------------------------------------------------------

This work is motivated by a clear gap: organizations increasingly need to
process PDFs programmatically, but the tools available today either fail on
complex pages or apply the same strategy everywhere regardless of whether it
fits.

\subsection{Growing Importance of Document Digitization}

Document digitization is accelerating across all sectors. Organizations are
not just scanning documents anymore. They need to extract, query, and
transform document content as part of larger automated workflows. This is
driven partly by the growth of AI-based systems that consume structured text,
such as LLM-based assistants and RAG pipelines, where the quality of the
input directly affects the quality of the output \autocite{khan2024}.
PDF-to-Markdown conversion is a key step in this pipeline: it bridges the
gap between static archived documents and the structured, searchable
knowledge bases that modern AI workflows depend on.

Scientific knowledge is a clear example. Blecher et al.\ \autocite{blecher2023}
note that PDFs make up 2.4\% of Common Crawl and are the second most
common data format on the internet after HTML. Most of that content is
not directly machine-readable. For AI systems that need to reason over large
document collections, this is a serious bottleneck. Khan et al.\ \autocite{khan2024}
tested this directly: a RAG system built on poorly extracted PDF text produced
low-quality answers even when the correct information was present in the
documents. The extraction step, not the language model, was the limiting factor.

The problem is not limited to academic papers. Ding et al.\ \autocite{ding2026}
identify finance, healthcare and legal domains as sectors where visually rich
documents are the primary carrier of information. In all three, automated
workflows increasingly need to read, compare and summarise document content
at scale. PDF was not designed for this. It was designed for printing.
Closing that gap between the printed format and machine-readable output is
what motivates the work in this thesis.

\subsection{Advantages of Markdown as a Structured Representation}

Among common output formats, Markdown is the most practical choice for the
use cases this thesis targets. It is human-readable without a renderer,
compatible with Git for version control, easy to parse programmatically, and
well-supported by LLM and RAG tooling.

The choice matters for downstream tasks. Fragmented or unstructured text
produces poor embeddings, which causes retrieval to fail even when the relevant
content is in the document. Markdown preserves hierarchy through headings, lists
and tables, which means the document's structure is visible to the embedding
model and to any downstream reader. A plain paragraph dump does not provide that.

Markdown also has a practical advantage over XML and HTML for this setting.
Both of those formats require a parser to read, and neither is easy to inspect
by hand. Markdown is plain text. A developer can open the output file in any
text editor and immediately see whether the conversion was correct. For a
system that will be evaluated by comparing its output to hand-written reference
files, this readability matters: it makes the reference creation process
faster and the error analysis more direct.

\subsection{Comparison of Extraction Strategies}

The four strategies evaluated in this thesis differ in how they handle
different content types. Text-only extraction is fast and does not require
an LLM, but it cannot process scanned pages or pages dominated by images
or formulas. Image-only extraction sends every page through a vision LLM,
which handles all content types but applies that cost even to plain-text
pages that do not need it. The hybrid strategy improves structural fidelity
by combining both text and image signals, but still applies the same pipeline
to every page uniformly. The adaptive approach routes each page to the
strategy that fits its detected content type, which is the key difference
none of the other three provide.

Meuschke et al.\ \autocite{meuschke2023} evaluated mainstream PDF parsing tools
across multiple document types and found that no single tool performed
consistently well across all of them. Rule-based text extractors performed best
on simple single-column documents but failed on scientific and patent documents.
This finding directly motivates the adaptive approach: if the best tool for one
document type is the wrong tool for another, a system that applies one strategy
uniformly will always underperform on at least some content.

Ingemarsson and Persson \autocite{ingemarsson2024} showed the cost side of the
same trade-off: multimodal tools scored highest overall but applied the same
high computation to every page regardless of content. The right extraction
method depends on what a specific page contains, not on the document as a
whole. Per-page adaptive routing is a direct answer to that observation.

% ----------------------------------------------------------------
\section{Research Objectives and Questions}
% ----------------------------------------------------------------

The main goal of this thesis is to design, implement, and evaluate a multimodal,
content-aware system that converts PDF documents into structured Markdown by
selecting an extraction strategy per page based on detected content type. The
system uses LM Studio as a local LLM runtime to ensure full reproducibility
without dependence on external cloud services. The central question is whether
this per-page selection outperforms any single fixed strategy on documents with
diverse content types, while reducing computational cost by skipping the LLM on
plain-text pages.

In practice, plain-text pages are routed to PyMuPDF and visually complex pages
are sent to a locally hosted vision LLM via LM Studio. A single document can
use both pathways, depending on what each page contains.

This leads to five research questions.

\textbf{RQ1 --- Adaptivity.} On documents with diverse content types, does
per-page strategy selection achieve higher accuracy and structural fidelity
than any single fixed strategy applied uniformly to the full document?

\textbf{RQ2 --- Strategy comparison.} How do text-only, image-only, and hybrid
strategies compare in terms of content and structural accuracy across different
bachelor and master thesis documents?

\textbf{RQ3 --- Efficiency.} What are the computational costs (processing time,
token usage, and LLM call count) of each strategy and of the adaptive system,
and how do those costs relate to accuracy?

\textbf{RQ4 --- Output quality.} How can Markdown quality be reliably measured
beyond BLEU (Bilingual Evaluation Understudy) and ROUGE (Recall-Oriented
Understudy for Gisting Evaluation), given that structured output requires
evaluating heading
hierarchy, list fidelity, and table reconstruction?

\textbf{RQ5 --- Consistency.} Do the performance differences between strategies
hold consistently across different bachelor and master thesis documents, or does
the relative ranking between strategies depend on the specific document?

These questions are tested against four hypotheses, which are evaluated against
the experimental results in Chapter~6.

\textbf{H1 --- Adaptive superiority.} On documents with diverse page content,
the adaptive system should outperform any single fixed strategy applied
uniformly, because it matches the best-suited strategy to each page type.
On documents with only one content type, adaptive should equal the
corresponding fixed strategy rather than exceed it.

\textbf{H2 --- Strategy-content fit.} Text-only will work well on plain-text
pages but fail on image-heavy or formula-rich content. Image-only and hybrid
strategies should handle diverse content well. Both treat every page as an
image and pass it through a vision LLM regardless of what the page actually
contains, which produces higher coverage but at significant computational cost.

\textbf{H3 --- Efficiency of the adaptive approach.} Despite the extra
classification step, the adaptive system should cost less on average than
image-only or hybrid. Plain-text pages skip the LLM entirely, and non-text
pages receive targeted prompts rather than a single general-purpose one,
which reduces unnecessary inference.

\textbf{H4 --- Hybrid as a strong baseline.} For documents with a balanced
mix of text and visual content, the hybrid strategy should be a competitive
baseline. The adaptive system is expected to match or beat it by avoiding
unnecessary LLM calls on plain-text pages.

% ----------------------------------------------------------------
\section{Scope and Delimitations}
% ----------------------------------------------------------------

This thesis works only with PDF as the input format. The system is tested
on bachelor and master thesis documents in German and English. The system
architecture is compatible with any LLM backend that supports a vision
Application Programming Interface (API),
but all experiments are conducted locally using LM Studio, so there is no
cloud dependency and results are fully reproducible.

The system detects the document language from the first page using the
\texttt{langdetect} library and adjusts LLM prompts to prevent the model
from translating content. If the page text is too short for reliable
detection, the system defaults to English. No language-specific training or
fine-tuning is applied, and the evaluation metrics are language-agnostic.

The following are out of scope: real-time or streaming document processing,
input formats other than PDF (e.g., DOCX, HTML, PPTX), fine-tuning or
task-specific training of LLM models, enterprise-scale deployment, cloud
infrastructure, and support for right-to-left scripts or ideographic writing
systems. The output format is Markdown only. Conversion to other structured
formats such as LaTeX, HTML, or JSON is not evaluated here, though the
architecture could support it through prompt changes.

% ----------------------------------------------------------------
\section{Thesis Contributions}
% ----------------------------------------------------------------

The central contribution is a per-page content classifier that assigns each
PDF page to one of six types: TEXT, IMAGE, FORMULA, TABLE, MIXED, or EMPTY.
It uses structural signals from PyMuPDF: embedded image count, vector path
density, presence of mathematical Unicode symbols, table detection, and
character count. No LLM is involved at this stage, so classification is fast
and deterministic.

For text pages, heading levels are inferred from font sizes relative to the
dominant body text size on the page. For pages routed to the vision path,
embedded figures are extracted using PyMuPDF's image extraction API and
saved as individual files. The vision LLM is instructed to include these paths
as Markdown image links, so figures are preserved in the output.

Output quality is measured across seven metrics: text similarity, heading
structure, list structure, table structure, code block score, paragraph
structure, and word overlap. These cover aspects of structural fidelity that
BLEU and ROUGE cannot capture. A configurable experiment runner tests all
combinations of strategy, model, prompt variant, temperature, and page across
the benchmark documents, records scores and token usage per page, and generates
comparison reports.

The full system runs locally using PyMuPDF, pymupdf4llm, and LM Studio. All
prompt templates, evaluation scripts, experiment configurations, and test
suites are included in the repository.

% ----------------------------------------------------------------
\section{Thesis Structure}
% ----------------------------------------------------------------

The thesis builds in three phases. The first establishes the foundations and
situates the work. Chapter~2 covers the technical background: how PDF stores
content internally, what Markdown is and why it suits this use case, how
modern LLMs process documents, and how the three main extraction strategies
differ. Chapter~3 surveys existing tools and systems, from traditional PDF
parsers to multimodal AI models, and identifies the gap this work addresses.

The second phase covers the research design and implementation. Chapter~4
describes the methodology: the document dataset, ground truth construction,
evaluation metrics, and experimental setup. Chapter~5 documents the system
architecture, from the per-page classifier and four conversion strategies to
the LM Studio integration and post-processing pipeline.

The third phase presents the findings and their interpretation. Chapter~6
reports the experimental results, with scores by strategy, model, and
document, and qualitative analysis of where each strategy succeeds or fails.
Chapter~7 interprets those results against the five research questions and
four hypotheses, and examines threats to validity. Chapter~8 closes with a
summary of contributions, limitations, and directions for future work.
